\documentclass[11pt,a4paper]{article}
\usepackage[margin=26mm]{geometry}
\usepackage{amsmath,amssymb,amsthm,hyperref}
\newtheorem{proposition}{Proposition}
\title{A panel obstruction in an omitted Mathieu geometry}
\author{An addition to the residue shadow calculations}
\date{}
\begin{document}\maketitle
The rank-three $M_{12}$ entry in Tables 1 and 7 of Jacobs and
Leemans~\cite{source} gives eighteen eligible geometries, divided into
seven satisfying (IP), one satisfying (WIP) alone, and ten failing
(WIP). Eligibility includes firmness, flag transitivity, residual
connectedness, residual weak primitivity, local 2-transitivity, and
geometric $(\mathrm{IP})_2$. Our earlier calculation on the twenty active
rank-three inputs in the supplied archive found seventeen eligible
geometries, with counts $7,1,9$. Three fail local 2-transitivity.

There is a separate potential omission. The retained complete subgroup
census contains twenty-one rank-three classes under automorphisms of
the specified acting group and permutations of types. The exact
comparison with every raw archive family identifies twenty of them;
class 18 of that census is absent. This index belongs to our census,
not to the historical numbering in the paper. We determine whether
this additional class supplies an eighteenth eligible geometry.

\begin{proposition}
Let $G=M_{12}$ and let $(H_1,H_2,H_3)$ be the marked subgroup family
in retained census class 18. Its orders are
\[
 |G|=95040,\qquad (|H_1|,|H_2|,|H_3|)=(24,660,660),
 \qquad |H_1\cap H_2\cap H_3|=6.
\]
The geometry satisfies all of the above eligibility conditions except
local 2-transitivity. It satisfies the full geometric intersection
property. Its panel of type 1 is the action of $A_5$ on the ten
two-element subsets of a five-element set. In particular, it cannot
be added to the locally 2-transitive population.
\end{proposition}
\begin{proof}
The entire marked family is retained in
\path{prior_representatives.g}. The unchanged exact residue procedure
from the common shadow paper checks every base residue, including
panels, and every remaining type. Flag transitivity makes these
representatives exhaustive for residues of the same type. The resulting
certificate verifies firmness, residual connectedness, residual weak
primitivity, and the geometric intersection property in ranks two and
three. These are exact finite group and coset calculations; no property
is inferred solely from the displayed orders.

The failure of local 2-transitivity has a smaller independent proof.
Put $B=H_1\cap H_2\cap H_3$ and $K=H_2\cap H_3$. A bijective
homomorphism $K\longrightarrow A_5$ maps $B$ to
\[
 D=\langle(2\,4\,3),(1\,5)(2\,3)\rangle.
\]
The program \path{identify_panel.g} checks this homomorphism and
that $D$ is exactly the stabilizer of $\{1,5\}$. Thus the right
$B$-cosets in $K$ are identified equivariantly with the ten two-element
subsets. This also gives a proof that is independent of the shadow
procedure: \path{panels.g} constructs those ten right cosets directly
and computes the action of their point stabilizer.

The stabilizer $D$ has three orbits on the two-element subsets. The
fixed subset is $\{1,5\}$. There are three subsets disjoint from it,
and six meeting it in one point. The displayed three-cycle permutes
the complementary three points, while the displayed involution
interchanges 1 and 5 and two of those three points. These generators
act transitively on each of the latter two sets. The subdegrees are
therefore $1,3,6$. Since the point stabilizer is not transitive on
the other nine points, the panel action is not 2-transitive. The
other two panels each have degree two. This proves the stated
failure without relying on a previously recorded eligibility flag.
\end{proof}

Consequently, relative to the retained complete census and its exact
archive correspondence, there are seventeen eligible rank-three
$M_{12}$ classes and the shadow counts remain $7,1,9$. This closes
the possible missing-input explanation of the earlier discrepancy.
It does not identify any class with a historical geometry number in
Tables 2 or 8, nor does it justify changing the stated equivalence
convention. The correspondence of historical labels remains a separate
obligation in a complete recovery of the paper.

The census, its subgroup catalogue and its completeness proof are
retained premises. The common residue theorem and its whole program
are also retained; they must be acquired in any comparison using this
application. The new calculation does not repeat the census. No
new original-paper recovery, open-conjecture result or Lean
formalization is claimed here. Software remains outside manuscript
length comparisons on both sides.

\begin{thebibliography}{9}
\bibitem{source} P.~Jacobs and D.~Leemans,
\emph{An algorithmic analysis of the intersection property},
LMS Journal of Computation and Mathematics \textbf{7} (2004), 284--299.
\bibitem{common} \emph{Shadows and residues in finite coset geometries},
the retained six-page common proof and its complete residue program.
\end{thebibliography}
\end{document}
